\documentclass[11pt]{article}

\usepackage[margin=1in]{geometry}
\usepackage{amsmath,amssymb}
\usepackage{enumitem}

\setlength{\parindent}{0pt}
\setlength{\parskip}{6pt}

\setlist[enumerate,1]{
    label=\textbf{\arabic*.},
    leftmargin=*,
    itemsep=1.1em
}
\setlist[enumerate,2]{
    label=(\alph*),
    leftmargin=2em,
    itemsep=0.3em,
    topsep=0.3em
}

\newcommand{\duedate}{Tuesday, September 8, by 11:59PM}

\begin{document}

\begin{center}
    {\Large\bfseries Math 465: Homework 1}\\[4pt]
    {\large Due: \duedate}
\end{center}

\hrule
\vspace{0.8em}

\textbf{Instructions.}
Write full and complete solutions to the following problems and submit them on Gradescope by the deadline. For each problem, explain the reasoning behind your answer; a formula or numerical answer without justification is not a complete solution.

At the beginning of your submission, acknowledge any permitted books, notes, or other sources you consulted, and list the names of any students with whom you discussed the homework.

\vspace{0.5em}
\hrule
\vspace{1em}

\begin{enumerate}

\item
A \emph{palindrome} is a string whose reversal is identical to the original string. How many strings of length \(n\) over the alphabet $\{0,1,2\}$ are palindromes?

A palindrome is determined only by half of characters. If n is odd, it includes middle one. so there are $3^{\lceil \frac{n}{2} \rceil}$ palindromes.

\item
How many ways are there to seat seven distinct people around a circular table if two seatings are considered the same whenever every person has the same two neighbors (regardless of whether they are left or right neighbors)?

First we fix one person, and other six people has $6!$ orders. However, because one order and its reversal count twice. So, there are $\frac{6!}{2}=360$ different ways.

\item
Eight distinct books are arranged in a row on a shelf. Two of the books are \emph{Pride and Prejudice} and \emph{Emma}.

\begin{enumerate}
    \item How many arrangements place these two books next to each other?

    First combine these two books together and regard them as a huge book. then there are 7 "books". Also, we need to consider the order of the two books. So there are $2\times7!=10080$ different ways.
    
    \item How many arrangements do not place these two books next to each other?

    If an order doesn't satisfy (a), then it satisfies (b). So there are $8!-10080=30240$ different ways.
    
    \item How many arrangements place \emph{Pride and Prejudice} somewhere to the left of \emph{Emma}?

    By symmetry the number of ways to put Pride and Prejudice to the left of Emma is the same as the way in different order. So there are $\frac{8!}{2}=20160$ ways.
    
\end{enumerate}

\item
How many six-digit positive integers have no repeated digits? 

All digit has 10 ways to choose, however the first digit cannot be 0, so it has 9 choices. We choose digit one by one. The first has 9 choice, the second has 10-1=9 choice, and so on. So there are $9\times9\times8\times7\times6\times5=136080$ numbers.

\item
Prove that the number of subsets of $[n]$ that have an odd number of elements is $2^{n-1}$.

Consider all the subsets of $[n-1]$, let $S$ is a subset in $[n-1]$, then consider $S$ and $S \cup \{n\}$, exactly one of them has an odd number of elements. There are $2^{n-1}$ of subsets in $[n-1]$. So there are $2^{n-1}$ subsets of $[n]$ that have an odd number of elements.

\item
A club has \(13\) undergraduate members and \(7\) graduate-student members. How many ways are there to form a committee of \(5\) members containing at least one undergraduate and at least one graduate student?

1 and 4: $\binom{13}{1} \binom{7}{4}$

2 and 3: $\binom{13}{2} \binom{7}{3}$

3 and 2: $\binom{13}{3} \binom{7}{2}$

4 and 1: $\binom{13}{4} \binom{7}{1}$

total: $\binom{13}{1} \binom{7}{4}+\binom{13}{2} \binom{7}{3}+\binom{13}{3} \binom{7}{2}+\binom{13}{4} \binom{7}{1}=14196$

\item
Consider the letters in the word
\[
\texttt{COMBINATORICS}.
\]

\begin{enumerate}
    \item How many distinct arrangements of these letters are there?

    It has 13 letters, with 2 C, 2 O,2 I.

    So there are $\frac{13!}{2!2!2!}=778377600$ permutations.
    
    \item How many distinct arrangements have the two copies of the letter \(\texttt{C}\) in nonconsecutive positions?

    Consider complement. If two Cs are consecutive, there are $\frac{12!}{2!2!}$ ways. So the complement has $778377600-\frac{12!}{2!2!}=658627200$ ways.
\end{enumerate}

\item
A class of ten students is divided into five two-person project teams. In how many ways can this be done?

We choose 2 students for 5 times. So it has $\binom{10}{2}\binom{8}{2}\binom{6}{2}\binom{4}{2}\binom{2}{2}$ ways.

However, we need to consider permutation of groups. For each division, there are $5!$ permutations. 

So there are $\frac{\binom{10}{2}\binom{8}{2}\binom{6}{2}\binom{4}{2}\binom{2}{2}}{5!}=945$ ways.

\item
A password is an eight-character string formed using uppercase letters and digits. There are \(26\) uppercase letters and \(10\) digits, and repetition is allowed. How many such passwords contain at least one digit?

Consider the complement. If all characters are letters, then it has $26^8$ ways. In total, it has $36^8$ ways. So there are $36^8-26^8=2612282842880$ passwords that contain at least one digit.

\item
Let \(k\geq 2\). How many strings of length \(n\) over an alphabet of \(k\) characters have no two consecutive characters equal?

Choose these n characters one by one. First character has k choices. After that, each character is different from the previous one, which has k-1 choices. So there are $k(k-1)^{n-1}$ strings that have no two conscutive characters equal.

\end{enumerate}

\end{document}